%% intro.tex


% hunsa's version of intro.tex
% Copyright (C) 2022

\chapter{Introduction to Scientific Writing}

In the following, we will give you a few hints on how to write your
thesis.

\section{Formatting Guidelines}

\subsection{Structure}
Structure your thesis properly. If you happen to have only one section
or subsection in a chapter or section, respectively, then remove the
heading.

For example, if you have
\begin{lstlisting}
Chapter 1 Introduction
1.1 Contribution
\end{lstlisting}
then change it into
\begin{lstlisting}
Chapter 1 Introduction
\end{lstlisting}

Similarly, if you have
\begin{lstlisting}
1.1 Quicksort
1.1.1 Randomized Quicksort
\end{lstlisting}
then, either you need a section with \verb|1.1.2| or remove the
heading \verb|1.1.1|.

\subsection{Capitalizing Headings}

All words (except ``a'', ``the'', ``in'', \etcet) are typically
capitalized (also conforms to guidelines of ACM and IEEE).

Example:
\begin{verbatim}
2 My Algorithm
2.1 Theoretical Foundation
2.2 NP-hardness Proof
2.3 Pseudocode of ABC Algorithm
2.4 Implementation using Java
2.5 Experimental Evaluation
\end{verbatim}

\section{Punctuation}

\subsection{Avoid Contractions}

Please do not use contractions in scientific writing, such as
``it's'', ``can't'', ``wanna'' or  ``won't''.

\subsection{Quotes}

Quoting seems to be very hard for some. When you want to refer to
someone else, you may use ``quotes''. But the quotes need to done
correctly, \ie, you start the quoted text with \verb|``| (two
backticks) and mark the end of the quote with \verb|''| (two
apostrophes). This ensures that the quotes are ``always set
correctly.''

\subsection{Footnotes}

You may want to use footnotes at times. However, if you do so, and
they happen to be at the end of the sentence, then put the period
before the footnote.\footnote{This is after the period.}

\subsection{Units}

The package \verb|siunitx| helps you with writing numbers and units.

For example, you may want to measure for \SI{100}{\micro\second} and
create file with \SI{10}{\gibi\byte} or \SI{10}{\giga\byte}, which
ever base you prefer.

\section{Floating Environments}

We usually try to align floats on top of the page, \eg,
\verb|{figure}[t]| or \verb|{table}[t]|. Sometimes it may be necessary
to modify this guidelines, \eg, on the first page, and then we can use
\verb|[h]|.

\subsection{Figures}

Images can be added with the \verb|\includegraphics| command as shown
in Figure~\ref{fig:intro}.  With \verb|\subcaption|, you can reference
subfigures, such as Figures~\ref{fig:intro:full width}
and~\ref{fig:intro:half width}.
%
\begin{figure}[h]
  \centering
  \begin{subfigure}[b]{0.45\columnwidth}
    \centering
    \includegraphics[width=\textwidth]{Logo-schwarz.pdf}
    \subcaption{The header logo at text width.}
    \label{fig:intro:full width}
  \end{subfigure}
  \begin{subfigure}[b]{0.45\columnwidth}
    \centering
    \includegraphics[width=0.5\textwidth]{Logo-schwarz.pdf}
    \subcaption{The header logo at half the text width.}
    \label{fig:intro:half width}
  \end{subfigure}
  \caption[Short caption for TOC]{The header logo at different sizes.} 
  \label{fig:intro} 
\end{figure}
%

Notice that the caption is below the figure. Please try avoiding pixel
graphics (jpeg, png) whenever possible.

\subsection{Tables}

\begin{table}[t]
  \centering
  \caption{Captions of tables are printed above.}
  \label{tab:example1}
  \begin{tabular}{ll}
    \toprule
    Thesis  & Template    \\
    \midrule
    Bachelor & \verb|example_hunsa.tex| \\
    Master  & \verb|example_hunsa.tex| \\
    PhD & \verb|example.tex| \\
    \bottomrule
  \end{tabular}
\end{table}

\tab~\ref{tab:example1} shows one example of inserting a table. Notice
that we print its caption above the table.


\subsection{Mathematical Expressions}

You can write mathematical expressions inline, such as
$\sum_{n=1}^{\infty} \frac{1}{n^2} = \frac{\pi^2}{6}$.
You may also write them in a block.
However, if you write an equation like the following one
\begin{equation*}
x = \sum_{i=1}^{n} i \quad ,
\end{equation*}
make sure that it is part of a sentence.

You can also write an equation, which you would like to reference later.
For example, the value of $y$ is defined as follows
\begin{equation}
\label{eq:ydef}
y = \sum_{i=1}^{n} 2^i \quad .
\end{equation}

Later, we may want to reference this equation. For example, our
definition of $y$ was given in Equation~\eqref{eq:ydef}.  Notice that
we have used \verb|\eqref| to reference this equation.

\subsection{Pseudo Codes and Listings}

Pseudo codes, algorithms, and other listings should also have a
caption above.

Let us give two example of how to work with algorithms and listings.

For writing algorithms, you can use the \verb|algorithmic| or
\verb|algorithmicx| package. Of course, there are also other packages
for achieving similar results. You may use other ones as well.

\begin{algorithm}[t]
\caption{Euclid’s algorithm}\label{alg:euclid}
\begin{algorithmic}[1]
  \Procedure{Euclid}{$a,b$}
  \State $r\gets a \mod b$
  \While{$r\not=0$}
  \State $a\gets b$
  \State $b\gets r$
  \State $r\gets a\bmod b$
  \EndWhile\label{euclidendwhile}
  \State \textbf{return} $b$
  \EndProcedure
\end{algorithmic}
\end{algorithm}

\alg~\ref{alg:euclid} shows an example algorithm, which was taken
directly from the documentation of the \verb|algorithmicx| package.

If you want to reference a particular line of the algorithm, you can
do that, but use \emph{Line}, \eg, in Line~\ref{euclidendwhile}, the
while loop is closed.

\begin{lstlisting}[language=C++,
  caption={Example of OpenMP program.},
  float=h,
  numbers=left]
#include <omp.h>

main ()  {

int nthreads, tid;

/* Fork a team of threads with each thread having a private tid variable */
#pragma omp parallel private(tid)
  {

  /* Obtain and print thread id */
  tid = omp_get_thread_num();
  printf("Hello World from thread = %d\n", tid);

  /* Only master thread does this */
  if (tid == 0) 
    {
    nthreads = omp_get_num_threads();
    printf("Number of threads = %d\n", nthreads);
    }

  }  /* All threads join master thread and terminate */

}
\end{lstlisting}

\section{Bibliography}

The bibliography is an important source of information for the reader.
Please ensure that each entry in the bibliography is completely
filled, \eg, each entry has a year and each conference paper has a
title of the proceedings it was published in.

Significant work on the decision problem has been done
before~\cite{Turing1936}.  Make sure that proper names and
abbreviations are surrounded with curly braces in your bib file, e.g.,
\verb|{MPI}| or \verb|{CUDA}|.

Turing~\cite{Turing1936} has done considerable work on the Entscheidungsproblem.

Terpstra~\etal~\cite{papi} showed how performance data can be collected with PAPI.

It could also happen that you will not need \etal, especially if there
are only two authors. For example, in 1972, Coffmann and
Graham~\cite{CoffmanG72} presented an optimal scheduling algorithm for
a system with two processors.

%%% Local Variables:
%%% mode: latex
%%% TeX-master: "example_hunsa"
%%% End:
